\documentclass[11pt]{article}
\usepackage{amsmath, amssymb}
\usepackage{hyperref}
\usepackage{graphicx}
\usepackage{tabto}

% Define \problemname and \illustration if not using Kattis template
\newcommand{\problemname}[1]{\section*{#1}}
\newcommand{\illustration}[3]{%
    \begin{figure}[h]
        \centering
        \includegraphics[width=#1\textwidth]{#2}
        \caption{#3}
    \end{figure}
}

\begin{document}

\problemname{Training Runs}

% Example use of the \illustration command
\illustration{0.7}{image.jpg}{Photo by \href{https://www.istockphoto.com/photos/hiking-trail-aerial}{iStock}}

A new trail running club is planning their weekly training plans. Each week, they have a specific number of runs planned and a circuit of trails available to them. They hope to \textbf{maximize total training load} each week, given that they can only run two segments per training session. They absolutely do not want to repeat a segment multiple times throughout the week; that would be so boring!

A run on segments i and i+1 has a training load of 2aiai+1-|ai-ai+1|: two demanding segments score high, but a sharp jump in difficulty within one run costs points.

The trail forms a linear sequence, where each segment connects to the segment that follows. The circuit is also not circular, so the last segment does not connect to the first segment.

Calculating the maximum total training load that can be accumulated each week is proving to be very time-consuming, since the circuits can get really big\dots Since they know you love this kind of puzzle, they have reached out to you to offer a program that will maximize the total training load in a given week.

Given the difficulty ratings of the N segments and the number of planned runs R, you have to select exactly R pairs of consecutive segments with no repeating segments, so that the total training load is as large as possible.

\section*{Input}

The input consists of one test case. The first line of the input consists of two integers, N and R: 
\begin{itemize}
\item N (2N5000), the number of segments in the trail circuit for that week.
\item R ($1 \leq R \leq \lfloor N/2 \rfloor$), the number of runs planned for that week.
\end{itemize}

The second line of input consists of N space-separated integers $a_1, a_2, \dots, a_N$ (0ai104), where ai represents the difficulty rating of segment i. The segments are ordered by how they connect sequentially along the trail.

\section*{Output}
A single integer representing \textbf{the maximum total training load with exactly R runs}, with each run covering two consecutive segments and no segment used more than once.

\section*{Sample Input 1    \tab Sample Output 1}
7 2       \tab 639 \newline
9 1 10 30 2 5 6 \newline

\end{document}

